\section{Prompts and Execution Traces}
\label{app:prompts-traces}

The official 119-topic RAG launcher used three LLM prompt templates in
the main runtime path: a combined sufficiency-and-query prompt, a dense
answer writer, and a sentence-level verify/revise pass. Citation
ordering and final structural validation are deterministic and use no
additional LLM prompt. The submission metadata
stores the base version \texttt{agentic\_rag\_baseline\_v1}; the writer
template identifies itself more specifically with the \texttt{-dense}
suffix. The templates below are transcribed from the runtime
implementation. Only typesetting has been normalized; runtime-injected
content is omitted.

The template text is reproduced verbatim; some heuristic instructions
and examples are not guaranteed to satisfy every stated constraint
simultaneously.

The combined sufficiency/query template is embedded in
the iteration runner, while the writer and verifier are shared
prompt-module templates. The offline checklist consumed by the sufficiency judge
(\S\ref{sec:rag-method}) is produced separately, once per topic, by a
single system/user decomposition call (retried on parse failure,
temperature 0). Stable source mappings, file hashes, and replay locations
are recorded in the project reproducibility record rather than exposed
as machine-specific paths in the paper.

\paragraph{Offline checklist generation (system and user prompt).} This
prompt runs once per topic, before the agentic loop, to decompose the
narrative into the sub-aspects the checklist enumerates. Its
sub-narrative decomposition produces the format consumed by the
sufficiency judge and dense writer. It does not use
organizer-provided nugget annotations or relevance judgments.

\begingroup
\footnotesize\ttfamily\raggedright
\textit{System:} You decompose a research narrative into the distinct
sub-aspects a complete answer must cover. You return strict JSON only -
no prose, no markdown fences.\\[4pt]
\textit{User:}\\
Narrative:\\
\mbox{[NARRATIVE]}\\[2pt]
Decompose this narrative into the distinct sub-aspects a complete,
well-rounded answer must cover. Return strict JSON with exactly this
shape:\\
\{"sub\_aspects": [\\
\phantom{\{}\{"title": "<short noun-phrase name of the aspect, 3-8 words>",\\
\phantom{\{}\ "importance": "vital" | "okay",\\
\phantom{\{}\ "expected\_facts": ["<one-sentence fact a good answer would state>", ...]\}\\
]\}\\[2pt]
Rules:\\
- 5 to 12 sub\_aspects, mutually distinct (no overlapping or duplicate
aspects).\\
- "vital" = the narrative explicitly asks about it or an answer
omitting it would clearly be incomplete; "okay" = secondary/supporting
aspect.\\
- 2 to 4 expected\_facts per aspect, each a single short declarative
sentence.\\
- Titles must be specific to this narrative, not generic ("Economic impact
of X on Y", not "Impacts").\\
- JSON only. No text before or after.
\par\endgroup

The resulting sub-aspects are mapped to checklist items by a
deterministic checklist-conversion step. The conversion rule is checked
against the 22-topic development checklist before applying it to the
119 test topics.

\paragraph{Evidence sufficiency and follow-up queries.}

\begingroup
\footnotesize\ttfamily\raggedright
Return ONLY one JSON object. No markdown. No explanation.\\
Required schema:
\{"enough":boolean,"queries":string[]\}\\
If the evidence already supports the topic, return\\
\{"enough":true,"queries":[]\}.\\
If evidence is insufficient, return\\
\{"enough":false,"queries":[...]\} with 1-3 BM25 keyword queries.\\
Query rules: 4-12 English tokens; keyword phrase, not sentence/question;
one aspect only; no duplicates.\\
Forbidden words: obtain, find, source, sources, detail, detailed,
comprehensive, concrete, examples, overview, history, impact, provide,
explain.\\
Previous queries: [PREVIOUS QUERIES]\\
Topic: [NARRATIVE]\\
Aspect checklist (use when judging sufficiency and choosing follow-up
queries; prefer queries that target aspects not yet covered by the
evidence):\\
\mbox{[CHECKLIST ITEMS]}\\
Evidence:\\
\mbox{[DOCUMENT IDS AND EXCERPTS]}
\par\endgroup

After an empty model response, the runtime varies the final retry line
to avoid replaying a cached empty response. The retry suffix does not
change the task definition.

\paragraph{Dense answer generation.}

\begingroup
\footnotesize\ttfamily\raggedright
Prompt version: agentic\_rag\_baseline\_v1-dense\\
You generate an evidence-grounded TREC RAG answer.\\
Use only the provided evidence documents. Do not use outside knowledge.\\[2pt]
LENGTH: the complete answer MUST total between 900 and 1010 words. Short
answers waste the official budget and lose coverage. Keep adding supported
factual sentences until you approach 1010 words.\\[2pt]
SENTENCE STYLE - every sentence must be:\\
- atomic: exactly one factual claim per sentence, at most
\textasciitilde35 words;\\
- concrete: preserve named entities, numbers, dates, organizations, places
verbatim from the evidence;\\
- grounded: cited to the evidence documents that state the fact.\\
- maximally SPECIFIC: prefer the most specific fact available - named
people, organizations, places, causes and their outcomes, findings and their
evidence. A precise claim (who did what, why, with what result) scores; a
general statement about the same aspect does not.\\
Do NOT write generic overview sentences that only describe the topic in the
abstract (e.g. 'This is an important and dynamic field', 'X is part of
culture'). Such sentences are worthless and must be omitted.\\
Good: "Persistent gender and racial pay gaps affect athlete compensation,
and business interests shape both media-rights deals and team management."\\
Bad: "Athlete compensation is a controversial and important topic."\\
No introduction, no conclusion, no hedging, no restating the question -
only supported facts.\\[2pt]
Organize the answer by these aspects of the topic; write roughly 80-120
words for EACH aspect (a few sentences each), in this order:\\
\mbox{[CHECKLIST ITEMS]}\\[2pt]
Citations are zero-indexed positions into the references array, not docids.\\
Each sentence may cite at most three references.\\
References must be unique ClimbMix docids and every reference must be cited
at least once.\\
Do not include uncited retrieved documents in references.\\
Return only strict JSON. Do not use Markdown fences, comments, or explanatory
prose.\\
Required JSON shape:\\
\{"references":["shard\_00000\_00000"],\\
"answer":[\{"text":"Supported sentence.","citations":[0]\}]\}\\
\mbox{[TOPIC ID AND NARRATIVE]}\\
Evidence documents:\\
\mbox{[UP TO 60 DOCUMENTS,}\\
\mbox{1,200 CHARACTERS EACH]}
\par\endgroup

\paragraph{Sentence-level verify/revise (production mode:
\texttt{weaken}).}

\begingroup
\footnotesize\ttfamily\raggedright
You are revising a cited RAG answer so that every sentence is fully supported
by its cited evidence.\\
For each sentence below, compare the sentence against its cited evidence
excerpts and act:\\
- Fully supported: keep the sentence unchanged with the same citations.\\
- Over-extended (partially supported): rewrite it so it claims ONLY what the
excerpts support. Do not add new claims.\\
- Not supported by its citations: rewrite it into a weaker related claim
that the excerpts DO support (hedge, narrow, or generalize it); if other
cited excerpts support the point, switch its citations to those; remove the
sentence only when nothing in any excerpt relates to it.\\
Rules: do not invent new references or citations; keep citation indices
pointing into the references array below; each sentence at most three
citations; keep the overall answer coherent and under 1024 words; keep the
original sentence order for kept/rewritten sentences.\\
Return only strict JSON, no Markdown fences, exactly this shape:\\
\{"references":["shard\_00000\_00000"],\\
"answer":[\{"text":"Supported sentence.","citations":[0]\}]\}\\
The references array must stay the same list (unused entries are allowed and
will be pruned automatically).\\
\mbox{[TOPIC ID AND NARRATIVE]}\\
References array (index -> docid): \mbox{[REFERENCES]}\\
Sentences to verify:\\
\mbox{[BATCH OF AT MOST 20 SENTENCES,}\\
\mbox{CITATIONS, AND EVIDENCE EXCERPTS]}
\par\endgroup

\paragraph{Citation ordering and final validation (no LLM prompt).}
After verify/revise, the runtime sends each answer sentence and its
cited document IDs to the evidence-scoring service with a 200-character
evidence budget. It sorts multiple citations by the returned evidence
score. If the service is unavailable, it leaves the original order
unchanged. Citation reordering therefore has no LLM prompt to
reproduce.

\paragraph{Worked trace: development topic 144.} The narrative asks for
a comparison of why banks such as Silicon Valley Bank fail, the hidden
risks to depositors, and how banks differ from credit unions. The
sufficiency judge ran for three iterations before stopping:

\begin{itemize}
\setlength\itemsep{2pt}
\item \textbf{Iteration 0} (12 documents read): \texttt{enough=false};
follow-up queries requested on FDIC/NCUA regulatory oversight,
credit-union community lending, and depositor/investor risk exposure.
\item \textbf{Iteration 1} (10 more documents read): \texttt{enough=false};
follow-up queries requested on ownership/governance differences and
service/fee comparisons.
\item \textbf{Iteration 2}: \texttt{enough=true}; loop stops with
\texttt{stop\_reason=enough} after 3 iterations total.
\end{itemize}

The summary is drawn from the complete stored development record for topic
144: its judge trace, iteration trace, and per-topic execution status
records. It is a development example and is not taken from
the separately archived official 119-topic run. Exact artifact identifiers
and hashes are provided in the project reproducibility record.

\paragraph{Stopping Behavior on the 119 Test Topics.} All 119 official topics
completed. Table~\ref{tab:trace-topology} aggregates the per-topic
execution status records: most topics exhaust the reading budget rather
than stopping because the judge reports sufficiency, and most reach the
6-iteration cap. The caps were active constraints, not fallback limits:
82 topics stopped at the document budget and 89 reached iteration 6.

\begin{center}
\captionof{table}{Stop reason and iteration count across all 119
official topics, aggregated from per-topic execution status records.}
\label{tab:trace-topology}
\centering
\small
\begin{tabular}{@{}lrr@{}}
\toprule
Stop reason & Topics & Share \\
\midrule
Maximum documents read (60) & 82 & 68.9\% \\
Judge reports sufficient (\texttt{enough}) & 36 & 30.3\% \\
No valid follow-up query & 1 & 0.8\% \\
\bottomrule
\end{tabular}

\smallskip

\begin{tabular}{@{}lr@{}}
\toprule
Iterations used & Topics \\
\midrule
2 & 5 \\
3 & 8 \\
4 & 4 \\
5 & 13 \\
6 (cap) & 89 \\
\bottomrule
\end{tabular}
\end{center}

Mean iterations were 5.45 (median 6, cap 6). Documents read
successfully ranged from 22 to 60 (mean 55.04, median 60, cap 60); no
topic failed outright. All 119 topics completed without runtime failures.
